\section*{Capítulo \ref{ch:g9:trig} --- Trigonometria no triângulo retângulo}

\begin{solution}{exo:g9:trig:1}
Hipotenusa: $\sqrt{9^2 + 12^2} = \sqrt{81 + 144} = \sqrt{225} = 15$.

Outro cateto: $\sqrt{17^2 - 8^2} = \sqrt{289 - 64} = \sqrt{225} = 15$.
\end{solution}

\begin{solution}{exo:g9:trig:2}
$7^2 + 24^2 = 49 + 576 = 625 = 25^2$: sim, é retângulo (recíproca de
Pitágoras), com o \omterm{def:g3:shapes:rightangle}{ângulo reto} oposto ao lado $25$.

$5^2 + 6^2 = 61 \neq 64 = 8^2$: não é retângulo.
\end{solution}

\begin{solution}{exo:g9:trig:3}
A hipotenusa é $[EF]$ (oposta ao \omterm{def:g3:shapes:rightangle}{ângulo reto} em $D$). O cateto oposto a
$\theta = \widehat E$ é $[DF]$; o cateto adjacente é $[DE]$. Logo,
\[
\cos\theta = \frac{DE}{EF}, \qquad
\sin\theta = \frac{DF}{EF}, \qquad
\tan\theta = \frac{DF}{DE}.
\]
\end{solution}

\begin{solution}{exo:g9:trig:4}
Cateto oposto: $10 \sin 28^\circ \approx 10 \times 0.469 = 4.7$.
Cateto adjacente: $10 \cos 28^\circ \approx 10 \times 0.883 = 8.8$.
(Conferindo: $4.7^2 + 8.8^2 \approx 22 + 77 \approx 100$.)
\end{solution}

\begin{solution}{exo:g9:trig:5}
$\tan\theta = \dfrac{4.5}{6} = 0.75$, então
$\theta = \tan^{-1}(0.75) \approx 37^\circ$.
\end{solution}

\begin{solution}{exo:g9:trig:6}
No \omterm{def:g6:shapes:triangles}{triângulo retângulo} formado pelo observador, pelo ponto do chão logo
abaixo do drone e pelo drone: a altitude de $120$ m é o cateto oposto
ao ângulo de elevação, e a distância horizontal $d$ é o cateto
adjacente. Então
\[
\tan 32^\circ = \frac{120}{d},
\qquad
d = \frac{120}{\tan 32^\circ} \approx \frac{120}{0.625} = 192
\text{ m}.
\]
\end{solution}

\begin{solution}{exo:g9:trig:7}
A subida ($1.2$ m) é o cateto oposto ao ângulo; o comprimento $L$ da
rampa é a hipotenusa:
\[
\sin 5^\circ = \frac{1.2}{L},
\qquad
L = \frac{1.2}{\sin 5^\circ} \approx \frac{1.2}{0.0872} \approx 13.8
\text{ m}.
\]
A rampa precisa ter pelo menos $13.8$ m de comprimento.
\end{solution}

\begin{solution}{exo:g9:trig:8}
De $\cos^2\theta + \sin^2\theta = 1$: $\sin^2\theta = 1 - 0.36 = 0.64$,
e $\sin\theta > 0$ para um ângulo agudo, então $\sin\theta = 0.8$.
Daí, $\tan\theta = \dfrac{\sin\theta}{\cos\theta} = \dfrac{0.8}{0.6} =
\dfrac43$.
\end{solution}

\begin{solution}{exo:g9:trig:9}
\emph{\omterm{def:g6:shapes:triangles}{Triângulo retângulo} \omterm{def:g6:shapes:triangles}{isósceles}} de catetos $1$: hipotenusa
$\sqrt{1 + 1} = \sqrt2$; cada ângulo agudo é $45^\circ$, e
\[
\cos 45^\circ = \frac{1}{\sqrt2} = \frac{\sqrt2}{2} = \sin 45^\circ,
\qquad
\tan 45^\circ = \frac11 = 1 .
\]

\emph{\omterm{def:g6:shapes:triangles}{Triângulo equilátero}} de lado $1$: a altura dele o divide em dois
\omterm{def:g6:shapes:triangles}{triângulos retângulos} de hipotenusa $1$, base $\frac12$ e altura
$h = \sqrt{1 - \frac14} = \frac{\sqrt3}{2}$ (Pitágoras). Os ângulos são
$60^\circ$ (na base) e $30^\circ$ (no topo). Lendo as razões:
\[
\cos 60^\circ = \frac{1/2}{1} = \frac12, \quad
\sin 60^\circ = \frac{\sqrt3/2}{1} = \frac{\sqrt3}{2}, \quad
\cos 30^\circ = \frac{\sqrt3}{2}, \quad
\sin 30^\circ = \frac12 .
\]
\end{solution}

\begin{solution}{pb:g9:trig:1}
\textbf{1.} A altura acima da altura dos olhos é o cateto oposto a
$35^\circ$, com cateto adjacente $60$ m:
$h = 60 \tan 35^\circ \approx 60 \times 0.700 = 42$ m.

\textbf{2.} $\tan\theta = 0.12$ dá
$\theta = \tan^{-1}(0.12) \approx 6.8^\circ$ --- íngreme para uma
estrada, e ainda assim um ângulo modesto. Uma estrada de $45^\circ$ tem
$\tan 45^\circ = 1$: uma inclinação de $100\,\%$. (``Cem por cento''
está longe da vertical --- um mal-entendido predileto nas pistas de
esqui.)

\textbf{3.} Num \omterm{def:g6:shapes:triangles}{triângulo retângulo} de ângulos agudos $\theta$ e
$90^\circ - \theta$, o cateto oposto a um é adjacente ao outro, e a
hipotenusa é comum. Então
$\sin(90^\circ - \theta) = \frac{\text{o oposto a ele}}{h} =
\frac{\text{adjacente a }\theta}{h} = \cos\theta$, e, simetricamente,
$\cos(90^\circ - \theta) = \sin\theta$: o \emph{co}sseno é o \omterm{def:g9:trig:ratios}{seno} do
\emph{co}mplemento.

\textbf{4.} Altura: $8 \sin 60^\circ = 8 \times \frac{\sqrt3}{2} =
4\sqrt3 \approx 6.93$ m. Pé: $8 \cos 60^\circ = 4$ m, exatamente.

\textbf{5.} $0.574^2 + 0.819^2 = 0.329 + 0.671 = 1.000$ (a menos do
\omterm{def:g6:decimals:rounding}{arredondamento}); e $\left(\frac12\right)^2 +
\left(\frac{\sqrt3}{2}\right)^2 = \frac14 + \frac34 = 1$. Qualquer par
de valores que falhe na identidade denuncia uma leitura errada ou uma
calculadora no modo errado --- um detector de erros de graça.

\textbf{6.} $\tan 40^\circ = \dfrac hx$ e
$\tan 30^\circ = \dfrac{h}{x + 200}$.

\textbf{7.} $x = \dfrac{h}{\tan 40^\circ}$ e
$x + 200 = \dfrac{h}{\tan 30^\circ}$; subtraindo,
$200 = h\left(\frac{1}{\tan 30^\circ} -
\frac{1}{\tan 40^\circ}\right)$, de onde sai a fórmula. Números:
$\frac{1}{0.577} \approx 1.732$, $\frac{1}{0.839} \approx 1.192$,
\omterm{ex:g1:subtraction:difference}{diferença} $0.540$: $h \approx \frac{200}{0.540} \approx 370$ m.

\textbf{8.} $x = \frac{h}{\tan 40^\circ} \approx \frac{370}{0.839}
\approx 441$ m. Conferindo a partir de $A$:
$\frac{h}{x + 200} \approx \frac{370}{641} \approx 0.577 =
\tan 30^\circ$: coerente.

\textbf{9.} O método de uma estação precisa da distância horizontal até
o ponto \emph{diretamente abaixo do cume} --- impossível de medir
através de um desfiladeiro ou dentro de uma montanha. O método de duas
estações a troca por uma distância que você mesmo consegue percorrer a
pé: os $200$ m entre os seus dois pontos de observação.

\textbf{10.} $h = \dfrac{200}{\sqrt3 - \frac{1}{\sqrt3}}
= \dfrac{200}{\frac{2}{\sqrt3}} = 100\sqrt3 \approx 173$ m --- o
segundo andar da Torre Eiffel, medido com duas visadas e uma caminhada.

\textbf{11.} Linha de visão $d$, \omterm{def:g3:shapes:shapes}{raios} $R$ (até o ponto de contato,
perpendicular à linha de visão) e $R + h$ (até o olho): Pitágoras dá
$d^2 + R^2 = (R + h)^2 = R^2 + 2Rh + h^2$, então
$d^2 = 2Rh + h^2 \approx 2Rh$ ($h$ é minúsculo diante de $R$). Olhos a
$1.80$ m: $d \approx \sqrt{2 \times 6.37 \times 10^6 \times 1.8}
\approx 4\,800$ m --- o horizonte no mar fica a menos de cinco
quilômetros.

\textbf{12.} $d \approx \sqrt{2 \times 6.37 \times 10^6 \times 324}
\approx 64$ km. Regra dos marinheiros: $3.6\sqrt{1.8} \approx 4.8$ km e
$3.6\sqrt{324} = 64.8$ km --- as duas batem.

\textbf{13.} $\tan\theta = \frac{2}{60}$ dá
$\theta \approx 1.9^\circ$. Os $0.5^\circ$ da Lua são cerca de um
\emph{quarto} de polegar: a Lua parece enorme e é escondida por uma
unha --- o olho é um transferidor ruim.

\textbf{14.} Escadas: $\tan^{-1}\frac{17}{29} \approx 30^\circ$ ---
escadas confortáveis sobem quase exatamente nos $30^\circ$ dos valores
especiais. A $60^\circ$ numa profundidade de $29$ cm, o degrau teria
$29\tan 60^\circ = 29\sqrt3 \approx 50$ cm de altura: uma escada de mão,
não uma escadaria.

\textbf{15.} $\tan 10^\circ \approx 0.18$; $\tan 45^\circ = 1$;
$\tan 80^\circ \approx 5.7$; $\tan 89^\circ \approx 57.3$. Quando
$\theta$ se aproxima de $90^\circ$, o cateto adjacente encolhe até
sumir enquanto o cateto oposto se mantém: o \omterm{def:g3:division:remainder}{quociente} cresce além de
qualquer limite. Em $90^\circ$, exatamente, não há triângulo nem valor
--- o gráfico da \omterm{def:g9:trig:ratios}{tangente} sobe por uma assíntota vertical, a primeira
de muitas no volume do ensino médio.
\end{solution}
